\documentclass[11pt,a4paper]{article}
\usepackage[utf8]{vietnam}
\usepackage{amsmath,amssymb}
\usepackage{graphicx}

\begin{document}

@@DOCUMENT_HEADER@@

@@SECTION_1_HEADER@@

\textbf{Câu 1.} Nếu hai đường tròn không cắt nhau thì số điểm chung của hai đường tròn là

@@TAB@@\textbf{A.} $1$.@@TAB@@\textbf{B.} $2$.@@TAB@@\textbf{C.} $3$.@@TAB@@\textbf{D.} $0$.

\textbf{Câu 2.} Cho hai đường tròn $(O_1)$ và $(O_2)$ tiếp xúc ngoài tại $A$ và một đường thẳng $\mathrm{d}$ tiếp xúc với $(O_1)$ và $(O_2)$ lần lượt tại $B, C$. Tam giác $ABC$ là

@@TAB@@\textbf{A.} Tam giác cân.@@TAB@@\textbf{B.} Tam giác đều.@@TAB@@\textbf{C.} Tam giác vuông.@@TAB@@\textbf{D.} Tam giác vuông cân.

\textbf{Câu 3.} Cho đoạn thẳng $OO'$ và điểm $A$ nằm trên đoạn $OO'$ sao cho $OA = 2O'A$. Vị trí tương đối của đường tròn tâm $O$ bán kính $OA$ và đường tròn tâm $O'$ bán kính $O'A$ là

@@TAB@@\textbf{A.} Nằm ngoài nhau.@@TAB@@\textbf{B.} Cắt nhau.@@TAB@@\textbf{C.} Tiếp xúc ngoài.@@TAB@@\textbf{D.} Tiếp xúc trong.

\textbf{Câu 4.} Cho hai đường tròn $(O), (O')$ cắt nhau tại $A, B$ trong đó $O' \in (O)$. Kẻ đường kính $O'C$ của đường tròn $(O)$. Chọn khẳng định \textbf{sai}.

@@TAB@@\textbf{A.} $AC = CB$.@@TAB@@\textbf{B.} $\widehat{CBO'} = 90^\circ$.

@@TAB@@\textbf{C.} $CA, CB$ là hai tiếp tuyến của $(O')$.@@TAB@@\textbf{D.} $CA, CB$ là hai cát tuyến của $(O')$.

\textbf{Câu 5.} Cho hai đường tròn tiếp xúc ngoài $(O; R)$ và $(O'; r)$ với $R > r$ và khoảng cách giữa hai tâm $OO' = \mathrm{d}$. Khi đó

@@TAB@@\textbf{A.} $\mathrm{d} = R - r$.@@TAB@@\textbf{B.} $\mathrm{d} > R + r$.@@TAB@@\textbf{C.} $R - r < \mathrm{d} < R + r$.@@TAB@@\textbf{D.} $\mathrm{d} = R + r$.

\textbf{Câu 6.} Cho đường tròn $(O_1; 3\text{ cm})$ tiếp xúc ngoài với đường tròn $(O_2; 1\text{ cm})$. Vẽ bán kính $O_1B$ và $O_2C$ song song với nhau cùng thuộc nửa mặt phẳng bờ $O_1O_2$. Gọi $D$ là giao điểm của $BC$ và $O_1O_2$. Tính số đo góc $\widehat{BAC}$ (với $A$ là tiếp điểm của $(O_1)$ và $(O_2)$).

@@TAB@@\textbf{A.} $90^\circ$.@@TAB@@\textbf{B.} $60^\circ$.@@TAB@@\textbf{C.} $100^\circ$.@@TAB@@\textbf{D.} $80^\circ$.

\textbf{Câu 7.} Cho hai đường tròn $(O; 10\text{ cm})$ và $(O'; 5\text{ cm})$ cắt nhau tại $A, B$. Tính độ dài đoạn thẳng $OO'$ biết $AB = 8\text{ cm}$ và $O, O'$ nằm cùng phía đối với $AB$ (làm tròn kết quả đến chữ số thập phân thứ nhất).

@@TAB@@\textbf{A.} $OO' \approx 6{,}5\text{ cm}$.@@TAB@@\textbf{B.} $OO' \approx 6{,}1\text{ cm}$.@@TAB@@\textbf{C.} $OO' \approx 6{,}3\text{ cm}$.@@TAB@@\textbf{D.} $OO' \approx 6{,}2\text{ cm}$.

\textbf{Câu 8.} Cho hai đường tròn $(O; 20\text{ cm})$ và $(O'; 15\text{ cm})$ cắt nhau tại $A, B$. Tính độ dài đoạn thẳng $OO'$ biết $AB = 24\text{ cm}$ và $O, O'$ nằm cùng phía đối với $AB$.

@@TAB@@\textbf{A.} $OO' = 7\text{ cm}$.@@TAB@@\textbf{B.} $OO' = 8\text{ cm}$.@@TAB@@\textbf{C.} $OO' = 9\text{ cm}$.@@TAB@@\textbf{D.} $OO' = 25\text{ cm}$.

\textbf{Câu 9.} Cho hai đường tròn $(O), (O')$ tiếp xúc ngoài tại $A$. Kẻ tiếp tuyến chung ngoài $MN$ với $M \in (O), N \in (O')$. Gọi $P, Q$ lần lượt là điểm đối xứng với $M, N$ qua đường nối tâm $OO'$. Khi đó tứ giác $MNPQ$ là hình gì?

@@TAB@@\textbf{A.} Hình thang cân.@@TAB@@\textbf{B.} Hình thang.@@TAB@@\textbf{C.} Hình thang vuông.@@TAB@@\textbf{D.} Hình bình hành.

\textbf{Câu 10.} Cho đường tròn $(O; 6\text{ cm})$ và $(O'; 2\text{ cm})$ cắt nhau tại $A, B$ sao cho $OA$ là tiếp tuyến của $(O')$. Độ dài dây chung $AB$ bằng

@@TAB@@\textbf{A.} $AB = 3\sqrt{10}\text{ cm}$.@@TAB@@\textbf{B.} $AB = \dfrac{6\sqrt{10}}{5}\text{ cm}$.

@@TAB@@\textbf{C.} $AB = \dfrac{3\sqrt{10}}{5}\text{ cm}$.@@TAB@@\textbf{D.} $AB = \dfrac{\sqrt{10}}{5}\text{ cm}$.

\textbf{Câu 11.} Cho các đường tròn $(O_1; 10\text{ cm})$, $(O_2; 15\text{ cm})$ và $(O_3; 15\text{ cm})$ tiếp xúc ngoài với nhau đôi một. Hai đường tròn $(O_2)$ và $(O_3)$ tiếp xúc nhau tại điểm $A$. Đường tròn $(O_1)$ tiếp xúc với $(O_2)$ và $(O_3)$ lần lượt tại $C$ và $B$. Khẳng định nào sau đây là đúng?

@@TAB@@\textbf{A.} $O_1A$ là tiếp tuyến chung của $(O_2)$ và $(O_3)$.@@TAB@@\textbf{B.} $O_1A = 25\text{ cm}$.

@@TAB@@\textbf{C.} $O_1A = 15\text{ cm}$.@@TAB@@\textbf{D.} Cả A và B đều đúng.

\textbf{Câu 12.} Cho hai đường tròn $(O), (O')$ tiếp xúc ngoài tại $A$. Kẻ các đường kính $AB$ của $(O)$ và $AC$ của $(O')$, gọi $DE$ là tiếp tuyến chung ngoài của hai đường tròn $(D \in (O), E \in (O'))$. Gọi $M$ là giao điểm của $BD$ và $CE$. Tính diện tích tứ giác $ADME$ biết $\widehat{DOA} = 60^\circ$ và $OA = 6\text{ cm}$.

@@TAB@@\textbf{A.} $12\sqrt{3}\text{ cm}^2$.@@TAB@@\textbf{B.} $12\text{ cm}^2$.@@TAB@@\textbf{C.} $16\text{ cm}^2$.@@TAB@@\textbf{D.} $24\text{ cm}^2$.

\textbf{Câu 13.} Cho hai đường tròn $(O), (O')$ tiếp xúc ngoài tại $A$. Kẻ các đường kính $AB$ của $(O)$ và $AC$ của $(O')$, gọi $DE$ là tiếp tuyến chung ngoài của hai đường tròn $(D \in (O), E \in (O'))$. Gọi $M$ là giao điểm của $BD$ và $CE$. Tính diện tích tứ giác $ADME$ biết $\widehat{DOA} = 60^\circ$ và $OA = 8\text{ cm}$.

@@TAB@@\textbf{A.} $12\sqrt{3}\text{ cm}^2$.@@TAB@@\textbf{B.} $\dfrac{64\sqrt{3}}{3}\text{ cm}^2$.@@TAB@@\textbf{C.} $\dfrac{32\sqrt{3}}{3}\text{ cm}^2$.@@TAB@@\textbf{D.} $36\text{ cm}^2$.

\textbf{Câu 14.} Cho các đường tròn $(O_1; 10\text{ cm})$, $(O_2; 15\text{ cm})$ và $(O_3; 15\text{ cm})$ tiếp xúc ngoài với nhau đôi một. Hai đường tròn $(O_2)$ và $(O_3)$ tiếp xúc nhau tại điểm $A$. Đường tròn $(O_1)$ tiếp xúc với $(O_2)$ và $(O_3)$ lần lượt tại $C$ và $B$. Tính diện tích tam giác $ABC$.

@@TAB@@\textbf{A.} $36\text{ cm}^2$.@@TAB@@\textbf{B.} $72\text{ cm}^2$.@@TAB@@\textbf{C.} $144\text{ cm}^2$.@@TAB@@\textbf{D.} $96\text{ cm}^2$.

\textbf{Câu 15.} Cho hai đường tròn $(O; R)$ và $(O'; r)$ với $R > r$ tiếp xúc ngoài tại $A$. Vẽ các bán kính $OB \parallel O'D$ với $B, D$ cùng thuộc nửa mặt phẳng bờ $OO'$. Đường thẳng $DB$ và $OO'$ cắt nhau tại $I$. Tiếp tuyến chung ngoài $GH$ của $(O)$ và $(O')$ với $G, H$ nằm ở nửa mặt phẳng bờ $OO'$ không chứa $B, D$. Tính $OI$ theo $R$ và $r$.

@@TAB@@\textbf{A.} $OI = \dfrac{R+r}{R-r}$.@@TAB@@\textbf{B.} $OI = \dfrac{R-r}{R+r}$.

@@TAB@@\textbf{C.} $OI = \dfrac{R(R-r)}{R+r}$.@@TAB@@\textbf{D.} $OI = \dfrac{R(R+r)}{R-r}$.

\vspace{0.4cm}
\begin{center}
\textbf{BẢNG ĐÁP ÁN TRẮC NGHIỆM}
\end{center}

@@ANSWER_KEY_TABLE@@

\vspace{0.4cm}
\begin{center}
\textbf{--------- HẾT ---------}
\end{center}

\end{document}
